% ==========================================================================
%  Thesis Section -- The shinyswarm R Package
%
%  Standalone section describing the reusable R package prototype that
%  generalises ShinySwarm's checkpoint mechanism. This is thesis-ready prose
%  (Discussion / Future Work material) describing WHAT the artifact is and
%  WHY it matters. It is deliberately separate from thesis-notes.tex, which
%  records HOW the library was built and is kept for reproducibility.
% ==========================================================================
\documentclass[11pt,a4paper]{article}
\usepackage[margin=2.5cm]{geometry}
\usepackage{booktabs,tabularx}
\usepackage{hyperref}
\usepackage[T1]{fontenc}
\usepackage{lmodern}
\usepackage{listings}
\usepackage{xcolor}
\lstset{
  basicstyle=\ttfamily\small,
  breaklines=true,
  columns=fullflexible,
  keepspaces=true,
  frame=single,
  framesep=4pt,
  xleftmargin=4pt
}

\title{The \texttt{shinyswarm} R Package:\\
  A Reusable Full-State Capture and Restore Library for Shiny}
\author{Iva Gunjaca\\Master of Software Engineering, University of Amsterdam}
\date{\today}

\begin{document}
\maketitle

% --------------------------------------------------------------------------
\section{Overview}
% --------------------------------------------------------------------------

To validate that ShinySwarm's checkpoint mechanism generalises beyond the
specific microservice architecture in which it was developed, a standalone R
package, \texttt{shinyswarm}, was implemented as a prototype. The package
distils the full-state capture-and-restore pattern --- previously implemented
individually inside each of the eight Shiny applications --- into two functions
that any Shiny application can use, with no dependency on Spring Boot, Redis,
Kafka, or Kubernetes.

The package is a proof of concept. Its purpose in this thesis is to demonstrate
the \emph{feasibility and generality} of the design proposed in the future-work
notes: that exact, reproducible state sharing for interactive Shiny applications
can be expressed as a small, transport-agnostic library. The prototype passes
\texttt{R~CMD~check} without errors, warnings, or notes, and is covered by an
automated test suite (Section~\ref{sec:testing}).

% --------------------------------------------------------------------------
\section{Motivation}
% --------------------------------------------------------------------------

Every application in ShinySwarm follows the same checkpoint pattern: capture the
session state (inputs plus computed reactive values), serialise it to JSON, and
later restore it by pushing the values back into the running application. In the
eight-application suite this logic was duplicated, with each application manually
building JSON payloads using application-specific field names.

Existing Shiny state-sharing mechanisms --- \texttt{shinyURL} and the built-in
bookmarking facility --- transfer only \emph{inputs}, forcing the receiving
session to recompute every output. For deterministic analyses this merely costs
time. For non-deterministic analyses, such as a Monte Carlo simulation executed
without \texttt{set.seed()}, recomputation produces \emph{different} results, so
input-only sharing cannot reproduce what the original user saw.

The \texttt{shinyswarm} package addresses both problems. It standardises the
duplicated logic into two function calls, and it captures computed reactive
values \emph{alongside} inputs, so the restored session reproduces the original
outputs exactly rather than recomputing them.

% --------------------------------------------------------------------------
\section{Design}
% --------------------------------------------------------------------------

\subsection{Transport agnosticism}

The central design decision is that the library produces and consumes JSON
strings only. It contains no networking, persistence, or broker logic. How a
checkpoint is transmitted or stored --- a REST request, a Kafka message, a Redis
entry, a local file, a database row --- is left entirely to the caller. The same
\texttt{capture\_state()} output is therefore usable in every ShinySwarm
architecture and in a plain single-user Shiny application running on a laptop.

\subsection{Public interface}

The package exposes two core functions and two typed wrappers:

\begin{lstlisting}[language=R]
capture_state(input, ..., max_size = NULL, compress = TRUE)
restore_state(session, state_json, ..., file_dir = tempdir(),
              allow_rds = FALSE, max_size = NULL)
state_file(path)   # wrap a file to embed
state_rds(obj)     # wrap an R object to serialise
\end{lstlisting}

Because R's reactive system exposes no mechanism to discover
\texttt{reactiveVal} or \texttt{reactiveValues} objects automatically, the
reactive values to capture or restore are registered explicitly by name through
\texttt{...}. This is a deliberate consequence of Shiny's design, not a
limitation of the library.

\subsection{Typed wrappers}

A Shiny application holds several kinds of state, each requiring a different
capture strategy. Rather than proliferating functions, the library uses
type-hinted wrappers that tell \texttt{capture\_state()} how to handle each
piece of state. The result is always a single JSON string; the wrappers control
what goes inside it.

\begin{center}
\begin{tabular}{lll}
\toprule
\textbf{State type} & \textbf{Example} & \textbf{Capture strategy} \\
\midrule
Scalar input        & \texttt{input\$num1 = 42}    & stored as-is \\
Reactive value      & \texttt{result()}             & accessor called, value stored \\
Reactive dataframe  & \texttt{shared\_df()}         & converted to a JSON array \\
File reference       & \texttt{summary.csv}          & \texttt{state\_file()}: read, gzip, base64 \\
Binary R object      & trained model                 & \texttt{state\_rds()}: serialise, gzip, base64 \\
\bottomrule
\end{tabular}
\end{center}

This generalises the Climate Anomaly checkpoint enrichment, which in the
microservice implementation embeds a processed CSV file into the checkpoint. In
the architecture that logic lives in the Spring Boot backend; in the library it
reduces to a single \texttt{state\_file()} call inside the application's save
handler, removing the need for the backend to mount shared volumes or parse CSV.

\subsection{Usage example}

With the library, an application's save and restore logic each reduce to a single
call:

\begin{lstlisting}[language=R]
library(shiny)
library(shinyswarm)

server <- function(input, output, session) {
  result <- reactiveVal(0)
  observeEvent(input$calculate, result(input$num1 + input$num2))
  output$result <- renderText(result())

  observeEvent(input$save_btn, {
    json <- capture_state(input, result = result)
    writeLines(json, "checkpoint.json")   # transport chosen by caller
  })

  observeEvent(input$restore_btn, {
    json <- paste(readLines("checkpoint.json"), collapse = "\n")
    restore_state(session, json, result = result)
  })
}
\end{lstlisting}

Inputs are restored through Shiny's low-level \texttt{session\$sendInputMessage()}
mechanism; reactive values are set directly through their accessors.

% --------------------------------------------------------------------------
\section{Reproducibility Guarantee}
% --------------------------------------------------------------------------

The package's defining property is exact reproduction of non-deterministic
results. Because computed reactive values are stored in the checkpoint rather
than recomputed on restore, a value produced by a stochastic process is
preserved bit-for-bit. To make this hold, the JSON serialisation is configured
to retain full numeric precision (\texttt{jsonlite::toJSON(..., digits = NA)});
the default four-decimal rounding would silently break the guarantee. This was
confirmed by a test in which a Monte Carlo mean computed without a fixed seed is
captured, the live value is then changed to a different draw, and a restore
recovers the original value exactly.

% --------------------------------------------------------------------------
\section{Security Considerations}
% --------------------------------------------------------------------------

Building the prototype surfaced a consideration that is absent from input-only
sharing approaches and that connects directly to the cross-user restore scenario
of RQ5: \textbf{a checkpoint is untrusted input}. When user B restores a
checkpoint authored by user A, the checkpoint must be treated as data from an
untrusted source. The restore path therefore defends against three classes of
attack.

\begin{table}[h]
\centering
\begin{tabularx}{\textwidth}{lX}
\toprule
\textbf{Attack} & \textbf{Defence} \\
\midrule
Path traversal via a crafted embedded filename
  & Filenames are reduced to a safe basename; path separators, drive letters,
    home-directory expansion, and \texttt{..} are rejected, so a restored file
    cannot escape its target directory. \\[4pt]
Arbitrary code execution via \texttt{unserialize()}
  & Restoring a serialised object can execute code while the object is
    reconstructed. \texttt{state\_rds} restoration is therefore disabled by
    default (\texttt{allow\_rds = FALSE}) and must be explicitly enabled for
    trusted checkpoints. \\[4pt]
Memory exhaustion (``decompression bomb'')
  & A small compressed payload can expand to many gigabytes. Decompression is
    performed as a bounded stream that aborts as soon as the running output size
    exceeds \texttt{max\_size}, so peak memory stays bounded regardless of the
    claimed decompressed size. \\
\bottomrule
\end{tabularx}
\caption{Security defences in \texttt{restore\_state()}.}
\end{table}

This is a research observation in its own right: full-state sharing trades the
recomputation cost of input-only sharing for a serialisation surface that must
be treated as a trust boundary. Any system that shares computed state across
users inherits this obligation.

% --------------------------------------------------------------------------
\section{Testing and Verification}
\label{sec:testing}
% --------------------------------------------------------------------------

The package is verified at two levels. Unit tests exercise the capture/restore
logic directly, including serialisation round-trips, the typed wrappers, the
exact-reproducibility property for non-deterministic values, and the three
security defences (path traversal, the \texttt{allow\_rds} gate, size limits, and
the decompression-bomb guard). An integration test runs the save/restore cycle
against Shiny's actual reactive engine using \texttt{shiny::testServer()},
confirming the library works with real \texttt{reactiveVal} and
\texttt{reactiveValues} objects rather than test doubles.

The package passes \texttt{R~CMD~check} with a clean result (no errors, warnings,
or notes), the standard correctness bar for an R package.

% --------------------------------------------------------------------------
\section{Limitations and Further Work}
% --------------------------------------------------------------------------

The prototype has known limitations, several of which are inherent to Shiny's
design rather than to the library:

\begin{itemize}
  \item \textbf{File-input widgets} cannot be restored through
    \texttt{sendInputMessage()}; the \texttt{state\_file()} wrapper embeds the
    uploaded content as a workaround.
  \item \textbf{Shiny modules} use namespaced input identifiers. Restoring across
    module session boundaries is not yet supported.
  \item \textbf{Explicit registration} of reactive values is required, because R
    provides no API to enumerate them.
  \item \textbf{\texttt{reactiveValues} keys cannot be deleted} in Shiny; on
    restore, keys absent from the checkpoint are cleared to \texttt{NULL} so
    their values do not leak, rather than being removed entirely.
\end{itemize}

Natural extensions include support for Shiny modules through recursive
capture/restore across session hierarchies, a quantitative micro-benchmark of
capture/restore overhead across the eight applications, and publication on CRAN
to make the reproducibility contribution available to the wider R~Shiny
community without any microservice infrastructure.

% --------------------------------------------------------------------------
\section{Contribution}
% --------------------------------------------------------------------------

The \texttt{shinyswarm} package shows that the reproducibility mechanism
developed for the ShinySwarm architecture is not tied to that architecture. The
same guarantee --- exact reproduction of computed results across users,
machines, and time --- can be delivered by a minimal, transport-agnostic library
usable by any Shiny developer. In doing so it both validates the design proposed
as future work and isolates a general security obligation for any system that
shares computed application state.

\end{document}
